\documentclass[11pt,a4paper]{article}

\usepackage{fontspec}
\usepackage[portuguese]{babel}
\usepackage[margin=2.2cm,top=2.4cm,bottom=2.4cm]{geometry}
\usepackage{xcolor}
\usepackage{titlesec}
\usepackage{array}
\usepackage[hidelinks]{hyperref}

\setmainfont{Helvetica Neue}

\definecolor{azul}{HTML}{1B4F8C}
\definecolor{cinza}{HTML}{555555}
\definecolor{verde}{HTML}{2E7D4F}
\definecolor{laranja}{HTML}{B25E12}
\definecolor{vermelho}{HTML}{9C2B2B}

\setlength{\parindent}{0pt}
\setlength{\parskip}{6pt}

\titleformat{\section}{\Large\bfseries\color{azul}}{\thesection.}{0.4em}{}
\titlespacing*{\section}{0pt}{18pt}{6pt}
\titleformat{\subsection}{\normalsize\bfseries}{\thesubsection}{0.6em}{}
\titlespacing*{\subsection}{0pt}{12pt}{2pt}

% \ideia{titulo}{complexidade}{stack sugerida}{resumo}
\newcommand{\complex}[1]{{\small\color{cinza}\textbf{Complexidade:} #1}}
\newcommand{\ods}[1]{{\small\color{verde}\textbf{ODS:} #1}}

\hypersetup{
  pdftitle={Ideias de Projeto - Laboratorio de Desenvolvimento de Software},
  pdfauthor={Squad UFOPA},
}

\begin{document}

\begin{center}
{\Huge\bfseries\color{azul} Ideias de Projeto}\\[4pt]
{\large Laboratório de Desenvolvimento de Software --- UFOPA}\\[2pt]
{\color{cinza} Banco de ideias para os 5 desafios do Product Design Studio Canvas --- Santarém-PA}\\[6pt]
{\small\color{cinza} Documento de apoio à decisão do squad --- 26 de agosto de 2026}
\end{center}

\vspace{6pt}
\hrule
\vspace{10pt}

\section*{Como usar este documento}
\addcontentsline{toc}{section}{Como usar este documento}

São \textbf{20 ideias}, quatro para cada um dos cinco desafios apresentados na aula de kickoff.
Cada ideia traz título, resumo do que o sistema faz, o núcleo funcional mínimo (o que precisa
existir no MVP) e uma classificação de complexidade. O critério de calibragem foi deliberado:
nada que se resolva com um CRUD de duas tabelas e nada que exija hardware embarcado,
treinamento de modelo de aprendizado de máquina ou convênio institucional prévio para funcionar.

\textbf{Restrições consideradas em todas as ideias:}

\begin{itemize}
  \item Squad de até 3 pessoas, com entrega da primeira versão em 09/10/2026 e versão final em 21/11/2026.
  \item Entregáveis obrigatórios: código-fonte versionado, documentação (UML: casos de uso, classes, estado, sequência), banco de dados modelado e artigo no formato SBC.
  \item Até 28/08/2026 é preciso ter tema escolhido, funcionalidades do MVP definidas e pelo menos 5 telas de protótipo de média fidelidade.
  \item Contexto local: conectividade instável fora do perímetro urbano, celulares de entrada, população acostumada a resolver as coisas por grupo de WhatsApp.
\end{itemize}

\textbf{Leitura da classificação de complexidade:}
\textit{Baixa-média} é entregável com folga e sobra tempo para polir a documentação e o artigo;
\textit{média} é o ponto ideal para três pessoas em três meses; \textit{média-alta} exige que o
squad corte escopo já na primeira semana e mantenha disciplina de entrega.

Ao final há uma tabela comparativa e uma seção de critérios de decisão.

\newpage

\section{Mobilidade Inteligente}
\ods{11 --- Cidades e Comunidades Sustentáveis \quad 13 --- Ação Contra a Mudança Global do Clima}

\emph{Como melhorar a mobilidade em Santarém com tecnologia, reduzindo os impactos do trânsito?}

\subsection{Cadê o Busão --- rastreamento colaborativo das linhas urbanas}
\complex{média}

Aplicativo em que o próprio passageiro informa em qual linha embarcou; enquanto a viagem dura, o
celular dele funciona como sensor de posição da linha, e todo mundo que espera no ponto vê onde o
ônibus está e a previsão de chegada. Resolve o problema sem depender de equipamento embarcado na
frota nem de convênio com a empresa concessionária --- que é justamente o que trava esse tipo de
projeto em cidades médias. \textbf{Núcleo do MVP:} cadastro de linhas e pontos, sessão de
``estou nesta linha'', mapa em tempo real, estimativa de chegada por média histórica e avaliação
da viagem. A literatura confirma que previsão de chegada e consulta de horários são as
funcionalidades mais valorizadas em aplicativos de transporte público, mas também alerta que esse
tipo de ferramenta tende a beneficiar quem já tem mais letramento digital --- o que rende um
requisito não funcional interessante para o artigo: modo econômico de dados e interface para
celular de entrada.

\subsection{Mototáxi Legal --- corrida com tarifa transparente e rastreio de segurança}
\complex{média}

Plataforma feita para as cooperativas de mototáxi que já operam na cidade, e não para competir com
elas. O condutor entra com cadastro validado (foto, placa, permissão), o passageiro chama por
proximidade e vê a tarifa calculada por faixa de distância segundo a tabela da própria cooperativa,
sem preço dinâmico. \textbf{Núcleo do MVP:} cadastro e validação de condutor, chamada por
geolocalização, cálculo de tarifa por faixa, botão de compartilhar corrida em andamento com um
contato de confiança e histórico de corridas por condutor. O ganho de escopo aqui é institucional:
formaliza um modal que carrega boa parte do deslocamento urbano em Santarém e hoje não tem registro
nenhum.

\subsection{Rota Segura --- mapa colaborativo de riscos no trajeto}
\complex{baixa-média}

Qualquer pessoa registra um ponto crítico no trajeto --- buraco, faixa apagada, semáforo quebrado,
ponto de ônibus sem abrigo, esquina sem iluminação --- com foto, categoria e localização. O sistema
agrupa as ocorrências por bairro e gera um painel público de prioridades, além de um relatório
exportável para o órgão de trânsito. \textbf{Núcleo do MVP:} denúncia georreferenciada com foto,
confirmação por outros usuários (``também vi isso aqui''), fila de moderação, mapa de calor e
painel por bairro. É a ideia de menor risco técnico da área --- e a que mais depende de o squad
fazer um bom trabalho de UX e de modelagem para não virar mais um formulário de reclamação.

\subsection{Embarque Ribeirinho --- horários e lotação de lanchas e balsas}
\complex{média}

Santarém se conecta a Alter do Chão, às comunidades do Tapajós e à margem oposta por embarcações
cujos horários circulam em cartaz de madeira e grupo de WhatsApp. O aplicativo centraliza rotas,
horários, preços e status real das embarcações, com confirmação colaborativa de saída e de lotação
por quem está no porto. \textbf{Núcleo do MVP:} catálogo de embarcações e rotas, agenda de partidas,
check-in colaborativo de ``saiu'' / ``lotado'' / ``atrasou'', alerta de alteração por seca ou cheia
do rio e funcionamento offline com sincronização posterior. É a ideia mais autenticamente local do
desafio 1, e a que tem menos concorrente pronto no mercado.

\newpage

\section{Causa Animal e Cultura}
\ods{15 --- Vida Terrestre}

\emph{Como reduzir o abandono de cães e gatos em Santarém, registrando cada um de forma escalável
e mudando a cultura da população no trato com esses animais?}

\subsection{RG Animal --- identidade digital com QR code na coleira}
\complex{média}

Cadastro único de cães e gatos da cidade: foto, tutor, endereço, vacinas, castração e histórico.
Cada animal recebe um código único impresso em plaqueta de coleira; quem encontrar o animal
escaneia com a câmera do celular e fala direto com o tutor, sem depender de leitor de microchip.
\textbf{Núcleo do MVP:} cadastro de tutor e animal, geração e impressão do código, página pública
de leitura do código, canal de contato preservando o telefone do tutor e lembrete anual de
atualização de dados. O embasamento é forte: identificação permanente triplica a chance de o animal
voltar para o tutor, e cães identificados apresentaram taxa de devolução de 71\% contra 39\% dos não
identificados. O gargalo conhecido dessas políticas não é a tecnologia de identificação e sim o
cadastro desatualizado --- por isso o lembrete anual não é enfeite, é requisito.

\subsection{Achei um Bicho --- mural georreferenciado de perdidos e encontrados}
\complex{baixa-média}

Publicação de ``perdi'' e ``achei'' com foto, ponto no mapa e características; o sistema notifica
automaticamente quem segue aquele bairro e sugere pareamentos prováveis entre um cadastro de perdido
e um de encontrado (mesmo porte, mesma cor, mesma região, datas próximas). \textbf{Núcleo do MVP:}
publicação com foto e geolocalização, busca por raio, notificação por área de interesse, pareamento
por atributos e marcação de caso resolvido. A escolha de notificar por bairro e não por cidade
inteira tem base: 70\% dos cães recuperados estavam a menos de 1,6 km de casa, e 42\% a menos de um
quarteirão --- alerta de vizinhança funciona melhor que divulgação ampla.

\subsection{Castra+ --- agendamento e triagem das campanhas de castração}
\complex{média}

Sistema para organizar a fila das campanhas de castração e do castramóvel. O tutor se inscreve, o
sistema classifica por critério social e por cobertura do bairro, emite senha com data e hora,
envia lembrete e registra o procedimento no histórico do animal. \textbf{Núcleo do MVP:} inscrição,
fila com regras de priorização, agenda por unidade e data, emissão de senha, registro de
comparecimento e painel de cobertura por bairro. É a ideia com a lógica de negócio mais rica do
desafio 2 --- priorização, capacidade por dia, remarcação, no-show --- o que torna a modelagem UML
naturalmente interessante em vez de burocrática.

\subsection{Adota Santarém --- adoção com acompanhamento pós-adoção obrigatório}
\complex{média}

Vitrine dos animais de ONGs e protetores independentes, com formulário de perfil do adotante, termo
digital de responsabilidade e --- o diferencial --- check-ins obrigatórios com foto em 30, 90 e 180
dias após a adoção, com alerta para o protetor quando o check-in não acontece. \textbf{Núcleo do
MVP:} catálogo com filtros, candidatura à adoção, fluxo de aprovação pelo protetor, termo digital
e agenda de acompanhamento. Plataformas de adoção existem em abundância; o que a literatura aponta
como lacuna é exatamente o acompanhamento posterior, que é onde o abandono acontece de novo.

\newpage

\section{Revitalização do Centro Histórico}
\ods{8 --- Trabalho Decente \quad 9 --- Indústria e Infraestrutura \quad 11 --- Cidades Sustentáveis}

\emph{Como tornar o centro de Santarém um lugar atrativo para moradia, comércio, turismo, cultura e
lazer, com uso de tecnologias inovadoras?}

\subsection{Santarém 1661 --- trilha histórica com realidade aumentada por geolocalização}
\complex{média-alta}

Roteiro a pé pelo centro: ao chegar em um ponto --- Catedral, Praça Barão de Santarém, orla, Casa da
Cultura, prédios do porto --- o aplicativo libera narração em áudio, foto histórica sobreposta à
imagem da câmera e um selo colecionável. A realidade aumentada aqui é a versão leve, ancorada em
GPS e bússola, sem reconstrução tridimensional nem rastreamento de superfície. \textbf{Núcleo do
MVP:} roteiro com pontos, detecção de chegada por geolocalização, conteúdo multimídia por ponto,
sobreposição simples na câmera e progresso do visitante. Estudos com esse desenho exato ---
aplicativo com AR e GPS aplicado à reativação de centros históricos --- registraram ganho mensurável
de aprendizado e de engajamento, inclusive em versões pensadas para crianças e adolescentes. É a
ideia mais vistosa para o pitch final e a que mais exige controle de escopo: comece por três pontos,
não por doze.

\subsection{Centro Vivo --- vitrine digital do pequeno comércio do centro}
\complex{baixa-média}

Catálogo dos comerciantes da região central: o que vendem, horário, formas de pagamento,
localização e fotos. Busca por categoria, filtro de ``aberto agora'' e selo para
estabelecimentos históricos da cidade. \textbf{Núcleo do MVP:} cadastro autogerido pelo lojista,
catálogo de produtos ou serviços, busca por categoria e proximidade, horário de funcionamento com
lógica de aberto/fechado e canal de contato direto. Combate a percepção de esvaziamento do centro
com a arma mais barata que existe: informação organizada.

\subsection{Agenda Aberta --- o que está acontecendo no centro hoje}
\complex{baixa-média}

Plataforma colaborativa de eventos da cidade --- shows na orla, feiras, exposições, Sairé, missas,
festivais, sessões de cinema, roda de capoeira --- com submissão pelo produtor, curadoria pela
equipe e notificação por interesse. \textbf{Núcleo do MVP:} submissão de evento, fila de aprovação,
calendário com filtros por categoria e data, eventos recorrentes e favoritos com lembrete. Parece
simples à primeira vista, mas as regras de negócio (recorrência, conflito de horário e local,
cancelamento, moderação) dão bom material de modelagem sem risco técnico.

\subsection{Fachada Viva --- inventário colaborativo do patrimônio construído}
\complex{média}

Moradores, estudantes e o poder público fotografam e classificam os imóveis do centro: estado de
conservação, uso atual, ocupado ou vazio, elementos arquitetônicos preservados. O sistema monta um
mapa de calor de degradação e mantém o histórico de cada imóvel ao longo do tempo. \textbf{Núcleo
do MVP:} ficha do imóvel com foto e localização, formulário de avaliação padronizada, versionamento
das avaliações, mapa de calor e exportação dos dados. Entre todas as ideias deste documento, é a que
mais facilmente vira artigo SBC de verdade --- porque produz um conjunto de dados original sobre a
cidade, e não só um sistema.

\newpage

\section{Indústria Criativa e Resíduos}
\ods{11 --- Cidades Sustentáveis \quad 12 --- Consumo e Produção Responsáveis}

\emph{Como criar, por meio da tecnologia, um ecossistema que use o resíduo sólido como insumo de
uma nova indústria baseada em reutilização criativa?}

\subsection{Cata Fácil --- ponte entre quem gera, quem cata e quem compra}
\complex{média}

Plataforma de três perfis: o gerador (casa, comércio, repartição) publica que tem material
disponível e a janela de horário; o catador ou a cooperativa vê os chamados próximos no mapa e
reserva; o ferro-velho ou a indústria recicladora recebe o material consolidado. \textbf{Núcleo do
MVP:} três perfis com permissões distintas, chamado de coleta georreferenciado, reserva exclusiva
para evitar disputa, confirmação de coleta com peso e histórico por perfil. Existe protótipo
publicado com exatamente esse desenho para Marabá-PA, além de casos documentados no Recife e na
Cidade do Cabo --- o que dá base sólida de comparação para o artigo. Vale registrar a ressalva
honesta que a literatura faz: o aplicativo sozinho não transforma a vida do catador; o que ele
melhora, de forma mensurável, é a logística. Prometer isso no pitch, e não redenção social.

\subsection{Ecoponto Mais --- onde descartar o que não vai na lixeira}
\complex{baixa-média}

Mapa dos pontos de entrega voluntária da cidade com o que cada um aceita --- óleo de cozinha,
pilha, eletrônico, entulho, vidro, pneu, medicamento vencido --- horário de funcionamento e rota
até lá. \textbf{Núcleo do MVP:} cadastro de pontos com tipos de material aceitos, busca por tipo de
resíduo, filtro por proximidade e horário, rota até o ponto e sugestão de novo ponto pelo usuário.
Escopo enxuto e banco de dados bem definido; boa escolha para um squad que quer folga para caprichar
na documentação e no artigo.

\subsection{Feira do Reuso --- mercado de material descartado para quem cria}
\complex{média}

Empresas, marcenarias, gráficas e pessoas anunciam resíduo aproveitável --- madeira de embalagem,
retalho de tecido, PET, pneu, tambor, sobra de MDF, lona de banner --- e artesãos, marceneiros e
projetos de extensão procuram por tipo, quantidade e localização. O sistema contabiliza quanto
material foi desviado do aterro. \textbf{Núcleo do MVP:} anúncio com tipo, quantidade e foto, busca
com filtros, manifestação de interesse, combinação de retirada e painel de material desviado. É a
leitura mais fiel do enunciado do desafio: resíduo entrando como insumo de uma cadeia criativa, e
não apenas coleta seletiva com outro nome.

\subsection{Recicla Pontos --- gamificação da entrega com recompensa no comércio local}
\complex{média}

O usuário entrega recicláveis num ponto parceiro, o atendente registra o peso pelo aplicativo, o
sistema converte em pontos e o usuário troca por desconto em estabelecimentos parceiros.
\textbf{Núcleo do MVP:} carteira de pontos com extrato, registro de entrega com peso pelo ponto
parceiro, catálogo de recompensas, resgate com código de validação e ranking por bairro ou escola.
Avaliações de campo de plataformas assim registram a recompensa como o principal motivador declarado
de adesão, acima da consciência ambiental --- e o sistema de pontos, saldo, expiração e resgate
oferece regras de negócio ricas para exercitar orientação a objetos de verdade.

\newpage

\section{Rota Livre --- Hack Your City}
\ods{14 --- Vida na Água \quad e outras demandas críticas locais}

\emph{Liberdade total para o squad indicar o desafio local que vai resolver e propor um sistema
inovador para resolvê-lo.}

\subsection{Olho no Tapajós --- monitoramento cidadão da qualidade da água}
\complex{média}

Ribeirinhos, banhistas, pescadores e escolas registram observações padronizadas sobre o corpo
d'água --- cor, odor, presença de lixo, espuma, peixe morto, turbidez aparente --- com foto e ponto
no mapa, formando uma série histórica por local. \textbf{Núcleo do MVP:} formulário padronizado de
observação, captura com foto e coordenada, funcionamento offline com sincronização posterior,
histórico por ponto de monitoramento e painel público com filtro por período. Há base científica
direta: descritores qualitativos coletados por cidadãos têm valor científico reconhecido, e a
principal barreira relatada em iniciativas amazônicas desse tipo é conectividade e distância
tecnológica dos monitores rurais --- daí o modo offline ser requisito de primeira classe, não
funcionalidade futura.

\subsection{Régua do Rio --- alerta comunitário de seca e cheia}
\complex{média}

Voluntários de cada comunidade lançam periodicamente a cota do rio no seu trecho; o sistema cruza
com o histórico e emite alerta quando o nível se aproxima de faixas críticas para navegação,
abastecimento de água, acesso à escola e ao posto de saúde. \textbf{Núcleo do MVP:} cadastro de
pontos de leitura, registro de cota com foto da régua, série histórica com gráfico, faixas de alerta
configuráveis por ponto e notificação para quem segue aquele trecho. O tema é urgente depois das
secas extremas recentes do Tapajós e do Amazonas, e é o tipo de sistema que uma banca reconhece de
imediato como socialmente relevante.

\subsection{Diário de Pesca --- caderno digital da pesca artesanal}
\complex{média}

O pescador registra a captura do dia --- espécie, quantidade, arte de pesca, local, tempo de viagem
--- e recebe de volta o próprio histórico organizado e o preço médio praticado na região; a colônia
enxerga o dado agregado e anônimo. \textbf{Núcleo do MVP:} registro rápido de captura pensado para
uso na beira do rio, catálogo de espécies com foto, histórico pessoal, preço médio por espécie e
painel agregado para a colônia. Modelo já validado na Amazônia brasileira: monitoramento comunitário
por smartphone produziu dados com precisão comparável à de órgão oficial e custo menor, e existe
uma iniciativa de referência de escala amazônica --- o ICTIO --- para citar como trabalho
relacionado no artigo.

\subsection{Praia Segura --- condição diária das praias do Tapajós}
\complex{baixa-média}

Painel diário das praias de Alter do Chão, Maracanã e Ponta de Pedras: nível do rio, faixa de areia
exposta, lotação estimada, presença de salva-vidas, aviso de queda de barranco ou corrente forte,
alimentado por operadores de turismo, barraqueiros e visitantes. \textbf{Núcleo do MVP:} ficha por
praia, atualização diária por colaborador verificado, relato de condição pelo visitante com foto,
histórico dos últimos dias e alerta de risco. Escopo delimitado, valor imediato para o turismo e
ligação natural com a ODS 14 --- uma boa escolha para um squad que quer entregar algo redondo e bem
acabado em vez de ambicioso e incompleto.

\newpage

\section*{Tabela comparativa}
\addcontentsline{toc}{section}{Tabela comparativa}

\renewcommand{\arraystretch}{1.25}
{\small
\begin{tabular}{@{}p{0.5cm}p{5.4cm}p{2.4cm}p{6.6cm}@{}}
\hline
\textbf{\#} & \textbf{Ideia} & \textbf{Complexidade} & \textbf{Principal risco do projeto} \\
\hline
1.1 & Cadê o Busão & média & Precisa de massa crítica de usuários para o rastreio funcionar \\
1.2 & Mototáxi Legal & média & Depende de adesão de uma cooperativa real para validar \\
1.3 & Rota Segura & baixa-média & Virar formulário de reclamação sem uso prático \\
1.4 & Embarque Ribeirinho & média & Coleta inicial dos horários reais é trabalho de campo \\
2.1 & RG Animal & média & Cadastro desatualizado esvazia o sistema \\
2.2 & Achei um Bicho & baixa-média & Baixo volume de casos durante a validação \\
2.3 & Castra+ & média & Regras de priorização precisam de interlocutor real \\
2.4 & Adota Santarém & média & Concorrência com grupos de rede social já consolidados \\
3.1 & Santarém 1661 & média-alta & AR consome tempo desproporcional; escopo precisa de corte \\
3.2 & Centro Vivo & baixa-média & Cadastro inicial dos lojistas é trabalho manual \\
3.3 & Agenda Aberta & baixa-média & Pouca inovação percebida no pitch \\
3.4 & Fachada Viva & média & Padronizar a avaliação exige critério técnico bem definido \\
4.1 & Cata Fácil & média & Três perfis distintos multiplicam telas e casos de uso \\
4.2 & Ecoponto Mais & baixa-média & Escopo pode ficar pequeno demais para a nota \\
4.3 & Feira do Reuso & média & Mercado de dois lados: precisa dos dois para demonstrar \\
4.4 & Recicla Pontos & média & Recompensa exige parceiro comercial para ser crível \\
5.1 & Olho no Tapajós & média & Modo offline é obrigatório e custa tempo de implementação \\
5.2 & Régua do Rio & média & Sem leitura voluntária constante, a série histórica não existe \\
5.3 & Diário de Pesca & média & Usabilidade em campo é o requisito mais difícil \\
5.4 & Praia Segura & baixa-média & Depende de colaborador local atualizando todo dia \\
\hline
\end{tabular}
}

\section*{Critérios para a escolha}
\addcontentsline{toc}{section}{Critérios para a escolha}

Antes de votar em uma ideia, vale passar cada finalista por cinco perguntas:

\begin{enumerate}
  \item \textbf{Existe um usuário real acessível em uma semana?} Cooperativa, ONG, colônia de pesca,
        associação de moradores, secretaria. Sem interlocutor, a fase de requisitos vira ficção e o
        artigo perde a parte de validação.
  \item \textbf{O MVP cabe em cinco telas?} Se não couber, o escopo já está grande demais para a
        entrega de 28/08.
  \item \textbf{O sistema tem valor com dez usuários?} Ideias que só funcionam com mil pessoas
        (rastreamento colaborativo, mercado de dois lados) são demonstráveis, mas precisam de um
        plano de demonstração desde já.
  \item \textbf{Dá para demonstrar rodando no pitch de 21/11?} O pilar 3 da avaliação é o sistema
        funcionando, não o slide.
  \item \textbf{Existe literatura para citar no artigo SBC?} Todas as ideias deste documento têm ---
        as referências ao final são o ponto de partida.
\end{enumerate}

\textbf{Se o squad quiser uma recomendação:} para equilíbrio entre viabilidade e impacto,
\textit{Embarque Ribeirinho} (1.4), \textit{RG Animal} (2.1) e \textit{Fachada Viva} (3.4) são as
três que combinam escopo controlável, originalidade local difícil de encontrar pronta no mercado e
material sobrando para o artigo. Se a prioridade for o pitch mais impressionante, \textit{Santarém
1661} (3.1) --- com a ressalva de cortar o escopo para três pontos históricos logo na primeira
semana.

\newpage

\section*{Referências consultadas}
\addcontentsline{toc}{section}{Referências consultadas}

{\small
Levantamento feito na base Consensus (busca em literatura científica revisada por pares), agosto de 2026.

\begin{enumerate}
\item \href{https://consensus.app/papers/details/fc75220c52a95957a4dec4a2aa6f8901/}{Smarter but more unequal transport? How socioeconomic and digital inequalities hinder adoption of mobility apps in the Global South} --- Dodel et al., 2024, \textit{Travel Behaviour and Society}, 22 citações.
\item \href{https://consensus.app/papers/details/eb464c9e48945f00a7f8196af006dcc3/}{Analysis of Mobile Applications Supporting Public Transport in Smart Cities on the Example of Poland} --- Kolodziej et al., 2025, \textit{European Research Studies Journal}.
\item \href{https://consensus.app/papers/details/c3732851551955b693e3ee79154b597c/}{Travelers' Attitudes Toward Mobile Application-Based Public Transport Services in Lahore} --- Javid et al., 2021, \textit{SAGE Open}, 28 citações.
\item \href{https://consensus.app/papers/details/9df5e3297a9b57feb281f7123b5ab1d4/}{Epidemiology of Dog and Cat Abandonment in Spain (2008--2013)} --- Fatjó et al., 2015, \textit{Animals}, 58 citações.
\item \href{https://consensus.app/papers/details/c21ff3ff3de454c197b78a6de373f05e/}{A New Web-Based Tool for RTO-Focused Animal Shelter Data Analysis} --- Kremer, 2021, \textit{Frontiers in Veterinary Science}, 6 citações.
\item \href{https://consensus.app/papers/details/3196b2e85ce05a5fa3f71092ed09de2e/}{Characterization of animals with microchips entering animal shelters} --- Lord et al., 2009, \textit{JAVMA}, 58 citações.
\item \href{https://consensus.app/papers/details/222a5dfd4803519ba57e63637534c6a9/}{Awareness of and engagement with microchipping services among UK dog owners} --- Siettou, 2026, \textit{Discover Animals}.
\item \href{https://consensus.app/papers/details/4445082e39ec52a58cd14679b1191134/}{Development and Evaluation of a Digital Pet Recovery and Adoption Management System for Alaminos City} --- Morales et al., 2026, \textit{IJRISS}.
\item \href{https://consensus.app/papers/details/f8af04c8a17b5dfca393a7dfe326fb20/}{Augmented reality mobile apps for cultural heritage reactivation} --- Hincapié et al., 2021, \textit{Computers and Electrical Engineering}, 48 citações.
\item \href{https://consensus.app/papers/details/832428ca4b86569aa0c4d5fd56ce51e1/}{An Architecture for Mobile Outdoors Augmented Reality for Cultural Heritage} --- Panou et al., 2018, \textit{ISPRS International Journal of Geo-Information}, 96 citações.
\item \href{https://consensus.app/papers/details/598df49fb88c5b0b99f888422b88f243/}{Augmented reality and serious game to engage the alpha generation in urban cultural heritage} --- Capecchi et al., 2024, \textit{Journal of Cultural Heritage}, 35 citações.
\item \href{https://consensus.app/papers/details/66a13f72561658b9b24c7aead237b69d/}{Selective Collection: App Prototype for Circular Solid Waste Management in the Amazon} --- Flores et al., 2025, \textit{Latinoware 2025}.
\item \href{https://consensus.app/papers/details/dbe5bc5261c75b9fa150ce5aa75f7ccb/}{The use of ICT in the informal recycling sector: The Brazilian case of Relix} --- Coelho et al., 2019, \textit{EJISDC}, 15 citações.
\item \href{https://consensus.app/papers/details/4f8a9289f7715672b6b16b70b4695575/}{Effect of the ICT-enabled reclaimer system on the informal waste recycling system in Cape Town: The Regenize model} --- Siwawa, 2024, \textit{EJISDC}, 10 citações.
\item \href{https://consensus.app/papers/details/483ea6b908945b98b7a59f489108400b/}{A Mobile Platform for Smart Recycling and Eco-Friendly Sustainable Environment} --- Banik et al., 2025, \textit{Journal of Social Innovation and Knowledge}.
\item \href{https://consensus.app/papers/details/672f44792cda59b78db8d50d381c8134/}{``The water is murky, the water is not moving'': qualitative water quality assessment by citizen scientists} --- Blanco-Ramírez et al., 2025, \textit{Frontiers in Water}.
\item \href{https://consensus.app/papers/details/ed026e7aa0f251648d9112fcc58da08d/}{Implementation of a Global Citizen Science App in Community-Based Water Monitoring: Lessons Learned from CrowdWater in Latin America} --- Blanco-Ramírez et al., 2025, \textit{Citizen Science: Theory and Practice}, 8 citações.
\item \href{https://consensus.app/papers/details/4f22fe32a8c858848defb1a13455337f/}{Success factors for citizen science projects in water quality monitoring} --- San Llorente Capdevila et al., 2020, \textit{Science of the Total Environment}, 119 citações.
\item \href{https://consensus.app/papers/details/87069eba9f5157519d8e21171f0a9468/}{Community-based monitoring of small-scale fisheries with digital devices in Brazilian Amazon} --- Oviedo et al., 2017, \textit{Fisheries Management and Ecology}, 22 citações.
\item \href{https://consensus.app/papers/details/fca3b921d3d0590a84ce80c78678333c/}{Monitoreo pesquero participativo en la Amazonía peruana: motivaciones e intereses de usuarios de ICTIO} --- Rodriguez Flores et al., 2022, \textit{Neotropical Hydrobiology and Aquatic Conservation}.
\item \href{https://consensus.app/papers/details/583a95605c14578598bc101119a60fb0/}{Engagement in a Citizen Science Project in the Amazon Basin} --- Eyng et al., 2022, \textit{Citizen Science: Theory and Practice}, 7 citações.
\end{enumerate}
}

\end{document}
